%%%%%%%%%%%%%%%%%%%%%%%%%%%%%%%%%%%%%%%%%%%%%%%%%%%%%%%%%%%%%%%%%%%%%%
% Lớp 10 | Phương pháp tọa độ trong mặt phẳng (𝑂𝑥𝑦) | Tọa độ của véctơ
% Dạng 0H9?1-6: Toán thực tế  dùng hệ  tọa độ
% Số câu: 6  (sinh tự động bởi tach_cau.py; từ nay thầy sửa trực tiếp file này)
%%%%%%%%%%%%%%%%%%%%%%%%%%%%%%%%%%%%%%%%%%%%%%%%%%%%%%%%%%%%%%%%%%%%%%

%%%%%%%%%%%%%%%% TRẮC NGHIỆM 4 PHƯƠNG ÁN %%%%%%%%%%%%%%%%

% nguồn: MaHoaNH_TracNghiem_Lop10
\begin{ex}%[0H9C1-6]%[ma:c3a81446]
	Trong hệ tọa độ $Oxy$, cho hai điểm $A(2;-3)$, $B(3;-4)$. Biết $M(x;y)$ trên trục hoành sao cho chu vi tam giác $AMB$ nhỏ nhất. Giá trị của $x$ nằm trong khoảng nào sau  đây?
	 \choice
		 {\True $(2;3)$}
		{$(3;4)$} 
		{$(1;2)$} 
		{$(0;1)$} 
	\loigiai{ 
	 \immini{
		Nhận xét: $A$, $B$ nằm cùng phía  đối với trục hoành.\\
		Gọi $M(x;0)$ là điểm  cần tìm và $A'(2;3)$ đối xứng với $A$ qua trục hoành.\\
		 $\bullet$ $ \overrightarrow{A'B}=(1;-7)$.\\
		Ta có $P=AM+MB+AB=MB+MA'+AB\Rightarrow P\ge  A'B+AB$\\
		$\Rightarrow P_{\min}=A'B+AB\\\Leftrightarrow A'$,  $M$, $B$ thẳng hàng.\\
		$\bullet$  $\overrightarrow{A'M}=(x-2;-3)$.\\
		Ba điểm $A'$, $M$,  $B$ thẳng hàng $\Leftrightarrow \overrightarrow{AM}$ cùng phương $\overrightarrow{A'B}$ $$\Leftrightarrow-7(x-2)=-3\cdot 1\Leftrightarrow -7x+14=-3\Leftrightarrow x=\dfrac{17}{7}.$$
		Vậy $x\in(2;3)$. 
		}{ 
		\begin{tikzpicture}[line join=round,line cap=round,>=stealth,scale=1]
				\path 
					 (2,-3) coordinate (A)
					(3,-4) coordinate (B) 
					(2,3)  coordinate (A')
					(0,0) coordinate (O) 
					 (1,0) coordinate (x)
					(intersection of O--x and A'--B) coordinate  (M);
				\foreach \x/\g  in {B/-45,A'/90,A/-90,M/45}\fill[black](\x) circle (1pt) +(\g:3mm) node {$\x$};
				\draw[->]  (-1,0)--(4,0) node[below left] {$x$};
				\draw[->] (0,-5)--(0,5)  node[below left] {$y$};
				\draw  (0,0) node [below left] {$O$};
				\foreach  \x in {1,2,3,4}
				 \draw[thin] (\x,1pt)--(\x,-1pt) node [below left] {$\x$};
				 \foreach \y in {-4,-3,-2,-1,1,2,3,4}
				\draw[thin]  (1pt,\y)--(-1pt,\y) node [left] {$\y$};
				\draw  (A)--(M)--(A')
					 (M)--(B)--(A);
				 \draw[dashed] (0,-3)-|(2,3)--(0,3);
			\end{tikzpicture}
		} 
	} 
\end{ex}


%%%%%%%%%%%%%%%% TRẢ LỜI NGẮN %%%%%%%%%%%%%%%%

% nguồn: MaHoaNH_TraLoiNgan_Lop10
\begin{bt}%[0H9H1-6]%[ma:572d7145]
	Cho ba điểm $A(-1; 1)$, $B(2; 1)$, $C(-1;-3)$. Khi  đó diện tích của tam giác $ABC$ là
	\shortans{6} 
	 \loigiai{
		Ta có $A B=\sqrt{3^2+0^2}=3$, $A C=\sqrt{0^2+(-4)^2}=4$, $\overrightarrow{B C}=(-3;-4)$, $B  C=\sqrt{3^2+4^2}=5$.\\
		Dễ  thấy $A B^2+A C^2=B C^2$ nên $\Delta ABC$ vuông tại $A$.\\
		Diện tích tam giác là $S_{\Delta ABC}=\dfrac{1}{2} AB \cdot AC=\dfrac{1}{2} \cdot  3 \cdot 4=6$.
	 }
\end{bt}

% nguồn: MaHoaNH_TraLoiNgan_Lop10
\begin{bt}%[0H9C1-6]%[ma:0f98ede0]
	Sự chuyển động của  một tàu thủy được thể hiện trên một mặt phẳng toạ độ như sau: Tàu khởi hành từ vị trí $A(1; 2)$ chuyển động thẳng đều với vận tốc (tính theo giờ) được biểu thị bởi vectơ $\overrightarrow{v}=(3; 4)$. Xác định tung độ vị trí của tàu (trên mặt phẳng toạ độ) tại thời điểm sau khi khởi hành $1{,}5$ giờ.
	\shortans{8} 
	\loigiai{ 
	 \immini{Gọi $B(x; y)$, $(y>0)$; $|\overrightarrow{v}|=\sqrt{3^2+4^2}=5$; $\overrightarrow{A B}=(x-1; y-2)$.\\
		Quãng  đường tàu thủy chạy được sau $1{,}5$ giờ là $1{,}5\cdot 5=7{,}5$.\\
		$AB=\sqrt{(x-1)^2+(y-2)^2}=7{,}5 
		 \Leftrightarrow (x-1)^2+(y-2)^2=7{,}5^2.
		 $\hfill (1)\\
		$\overrightarrow{A B}$ và $\overrightarrow{v}$ cùng phương nên $\dfrac{x-1}{3}=\dfrac{y-2}{4} \Leftrightarrow x=\dfrac{3}{4} y-\dfrac{1}{2}$.  \hfill (2)\\
		Thay (2) vào (1) ta có }{\begin{tikzpicture}[line cap=round,line join=round,font=\footnotesize,>=stealth,scale=0.4]
		\draw[->] (-1,0)--(7,0)  node[above] {$x$};
		\draw[->]  (0,-1)--(0,9) node[left] {$y$};
		\clip (-1,-1) rectangle  (7,9);
		\path 
			(0,0)  coordinate (O)
			(1,2) coordinate (A) 
			({11/2},8) coordinate (B) 
			 (3,1) coordinate (V1)
			(6,5)  coordinate (V2);
		 \draw[dashed] (A)--(B)--($(B)+(3,4)$);
		\draw[dashed] (1,0)|-(0,2); 
		 \draw[->] (V1)--(V2) node[midway,below right] {$\overrightarrow{v}$};
		 \path (1,-0.4) node {$1$};
		\path  (-0.4,2) node {$2$};
		\foreach  \x/\g in {B/135,A/120,O/-135}{
			\fill (\x) circle (1pt) node[shift=(\g:3mm)]  {$\x$};}
	\end{tikzpicture}
		} 
		\[\left(\dfrac{3}{4} y-\dfrac{1}{2}-1\right)^2+(y-2)^2=7{,}5^2\Leftrightarrow 25 y^2-100 y-800=0  \Leftrightarrow \hoac{ &
		y=8  \Rightarrow x=\dfrac{11}{2} \\ & y=-4 \text{ (loại)}.}\]
	Vậy $B\left(\dfrac{11}{2};  8\right)$. }
\end{bt}


%%%%%%%%%%%%%%%% TỰ LUẬN %%%%%%%%%%%%%%%%

% nguồn: MaHoaNH_TuLuan_Lop10
\begin{bt}%[0H9H1-6]%[ma:b279dbb9]
	Cho $\Delta AB C$ có $A(5; 6)$, $B(4;-1)$, $C(-4; 3)$. Tọa độ điểm $M(a;b)$ thuộc đoạn  $BC$ sao cho $S_{\Delta MAB}=5 S_{\Delta MAC}$. Tính giá trị của biểu thức $a+b$.
	\loigiai{ 
		\begin{center}
			\begin{tikzpicture}[scale=1,font=\footnotesize,line join=round,line cap=round,>=stealth]
					 \def\a{4}
					\coordinate[label=left:\scriptsize $B$] (B)  at (0,0);
					\coordinate[label=right:\scriptsize  $C$] (C) at (\a,0);
					\coordinate[label=above:\scriptsize $A$] (A) at ($(B)+(65:.75*\a)$); 
					\coordinate[label=below:\scriptsize  $H$] (H) at ($(B)!(A)!(C)$);
					\coordinate[label=below:\scriptsize $M$] (M) at  ($(C)!1/6!(B)$);
					\draw  (M)--(A)--(B)--(C)--(A)--(H);
					\foreach \diem in  {A,B,C,M,H}\fill (\diem) circle (1.5pt);
					\newcommand{\gocv}[4][black]{\draw[#1]  ($(#3)!5pt!(#2)$)--($(#3)!2!($($(#3)!5pt!(#2)$)!.5!($(#3)!5pt!(#4)$)$)$)--($(#3)!5pt!(#4)$);}
					 \gocv{A}{H}{C}
				\end{tikzpicture}
		\end{center}
		Kẻ $AH \perp BC$. $S_{\Delta MAB}=5 S_{\Delta MAC}\Leftrightarrow \dfrac{1}{2}AH\cdot MB=5\cdot \dfrac{1}{2}AH\cdot MC\Leftrightarrow MB=5MC$. \\Mà $\overrightarrow{M B}$ và $\overrightarrow{M C}$ ngược hướng $\Rightarrow  \overrightarrow{M B}=-5 \overrightarrow{M C}$ \\ $\Leftrightarrow \heva{&x_B-x_M=-5(x_C-x_M)\\&y_B-y_M=-5(y_C-y_M)}\Leftrightarrow\heva{& 4-x_M=-5(-4-x_M)\\&-1-y_M=-5(3-y_M)}\Rightarrow M\left(-\dfrac{8}{3};\dfrac{7}{3}\right)$.\\
		Suy  ra $a+b=-\dfrac{1}{3}$.
	} 
\end{bt}

% nguồn: MaHoaNH_DungSai_Lop10
\begin{ex}%[0H9V1-6]%[ma:7b5f744f]
	Trên màn hình ra đa của đài kiểm soát không lưu (được coi như mặt phẳng tọa độ $ Oxy $ với đơn vị trên các trục tính theo ki-lô-mét), một máy bay trực thăng chuyển động thẳng đều từ thành phố $ A $ có tọa độ $ (600;200) $ đến thành phố $ B $ có tọa độ $ (200;500) $ và thời gian bay quãng đường $ AB $ là $ 3 $ giờ. Tọa độ $M(x_m,y_m)$ của máy bay trực thăng tại thời điểm sau khi xuất phát  $ 1 $ giờ thỏa mãn:
	\begin{choiceTF}
		{\True $x_m=\dfrac{1400}{3}$} 
		{\True  $y_m=300$}
		{\True $y_m > 200$ } 
		{$x_m  > 600$ }
	\end{choiceTF}
	 \loigiai{
		Giả sử $ M(x;y) $  là vị trí của máy bay trực thăng tại thời điểm sau xuất phát $ 1 $ giờ.\\
		Ta có $ \overrightarrow{AM}=(x-600;y-200) $, $  \overrightarrow{AB}=(-400;300) $.\\
		Vì máy bay trực thăng chuyển động thẳng đều nên $ \overrightarrow{AM}=\dfrac{1}{3}\overrightarrow{AB} $. Do  đó
		$$ \heva{&x-600=-\dfrac{400}{3}\\&y-200=100} \Leftrightarrow \heva{&x=\dfrac{1400}{3}\\&y=300.} $$ 
		Vậy vị trí của máy bay trực thăng tại thời điểm sau khi xuất phát $ 1 $ giờ là $ M\left(  \dfrac{1400}{3};300\right) $.
	} 
\end{ex}

% nguồn: MaHoaNH_DungSai_Lop10
\begin{ex}%[0H9V1-6]%[ma:10022c53]
	Để kéo đường dây điện băng qua một hồ hình chữ nhật $ABCD$ với độ dài $AB=200$ m, $AD=180$ m, người ta dự định làm $4$ cột điện liên tiếp cách đều, cột thứ nhất nằm trên bờ $AB$ và cách đỉnh $A$ khoảng cách $20$ m, cột thứ tư nằm trên bờ $CD$ và cách đỉnh  $C$ khoảng cách $30$ m. Gọi $\mathrm{d}(C_2,AB)$, $\mathrm{d}(C_2,AD)$ lần lượt là khoảng cách từ vị trí cột thứ hai đến các bờ $AB$, $AD$. Xác định tính đúng sai của các mệnh đề sau
	\begin{choiceTF}
		{\True $\mathrm{d}(C_2,AB)=60$  m}
		{$\mathrm{d}(C_2,AB)= 80$  m}
		{$\mathrm{d}(C_2,AB)=120$  m}
		{\True $\mathrm{d}(C_2,AD)=70$ m} 
	\end{choiceTF}
	 \loigiai{
	 \immini{
		Chọn hệ  tọa độ $Oxy$ sao cho $A(0;0)$, $B(200;0)$, $C(200;180)$, $D(0;180)$.\\
		 Gọi vị trí các cột điện được trồng là $C_1$, $C_2$, $C_3$, $C_4$.\\
		Do  $C_1$ thuộc cạnh $AB$ và $AC_1=20$ nên $C_1(20;0)$, do $C_4$ thuộc cạnh $CD$ và $C_4C=30$ nên $C_4(170;180)$.\\
		 Suy ra $\overrightarrow{C_1C_4}=(150;180)$. $(1)$\\
		Do bốn cột điện $C_1$, $C_2$, $C_3$, $C_4$ được trồng liên tiếp, cách đều trên một đường thẳng, nên $$\overrightarrow{C_1C_2}=\dfrac{1}{3}\overrightarrow{C_1C_4}\text{ và }  \overrightarrow{C_1C_3}=\dfrac{2}{3}\overrightarrow{C_1C_4}.$$
		Gọi tọa độ của $C_2$ đối với hệ trục đang xét  là $(x;y)$. Khi đó $\overrightarrow{C_1C_2}=(x-20;y)$.
		}{ 
		\begin{tikzpicture}[scale=0.75,font=\footnotesize,line join=round,line cap=round,>=stealth]
				\draw[->] (-1,0)--(11,0)  node[above]{$x$};
				\draw[->]  (0,-1)--(0,10) node[left]{$y$};
				\clip (-1,-1)  rectangle (11,10);
				 \draw[opacity=0.1] (-1,-1) grid (11,10);
				 \path
					(0,0) coordinate (O) 
					(0,0) coordinate  (A)
					(10,0) coordinate  (B)
					 (10,9) coordinate (C)
					(0,9) coordinate  (D)
					(1,0)  coordinate (C_1)
					 (8.5,9) coordinate (C_4)
					 ($(C_1)!{1/3}!(C_4)$) coordinate (C_2)
					($(C_1)!{2/3}!(C_4)$) coordinate (C_3); 
				 \draw (B)--(C)--(D)
					(C_1)--(C_4) 
					(3.5,0)|-(0,3) 
					 (6,0)|-(0,6);
				 \draw[fill=black] (10,0) circle (1pt) node[below]{$200$};
				\draw[fill=black] (0,9) circle  (1pt) node[left]{$180$};
				\foreach \diem/\g in {O/-135,A/-45,B/45,C/45,D/45,C_1/125,C_2/125,C_3/125,C_4/90} \draw[fill=black] (\diem) circle  (1pt)++(\g:0.4) node{$\diem$};
				\path (C_1)--(C_2)  node[midway,sloped,scale=0.8]{$||$};
				\path (C_2)--(C_3)  node[midway,sloped,scale=0.8]{$||$};
				\path (C_3)--(C_4)  node[midway,sloped,scale=0.8]{$||$};
			\end{tikzpicture}
		} 
		\noindent 
		Từ đó  và $(1)$, do $\overrightarrow{C_1C_2}=\dfrac{1}{3}\overrightarrow{C_1C_4}$ nên $\heva{&x-20=\dfrac{150}{3}=50\\&y=\dfrac{180}{3}=60.}$\\
		Suy ra $x=70$, $y=60$,  tức là $C_2(70;60)$.\\
		Khi  đó $\mathrm{d}(C_2,AB)=\mathrm{d}(C_2,Ox)=60$ (m) và $\mathrm{d}(C_2,AD)=\mathrm{d}(C_2,Oy)=70$ (m).
	 }
\end{ex}
